\documentclass[11pt]{article}
\usepackage[a4paper,margin=1.1in]{geometry}
\usepackage[T1]{fontenc}
\usepackage{lmodern}
\usepackage{amsmath,amssymb,amsthm}
\usepackage{booktabs}
\usepackage{longtable}
\usepackage{array}
\usepackage{microtype}
\usepackage[hidelinks]{hyperref}
\usepackage{natbib}
\usepackage{setspace}
\usepackage{enumitem}
\usepackage{xcolor}

\newtheorem{theorem}{Theorem}[section]
\newtheorem{proposition}[theorem]{Proposition}
\newtheorem{lemma}[theorem]{Lemma}
\newtheorem{corollary}[theorem]{Corollary}
\theoremstyle{definition}
\newtheorem{definition}[theorem]{Definition}
\newtheorem{assumption}[theorem]{Assumption}
\theoremstyle{remark}
\newtheorem{remark}[theorem]{Remark}

% A formal name, as it appears in the Lean development.
\newcommand{\leanunderscore}{\char`\_\hskip 0pt plus 0.4pt\relax}
\newcommand{\lean}[1]{\begingroup\ttfamily\small\let\_\leanunderscore #1\endgroup}
% The Lean name attached to a stated result.
\newcommand{\formal}[1]{\par\noindent\begingroup\footnotesize\raggedright\emph{Formal statement:} \lean{#1}.\par\endgroup}
\newcommand{\Thorn}{\text{\th}}
\newcommand{\R}{\mathbb{R}}
\newcommand{\E}{\mathbb{E}}
\DeclareMathOperator{\supp}{supp}

\onehalfspacing
\tolerance=600 \emergencystretch=1.5em

\title{Means-Tested Pensions and Monotone Policies:\\ Saving Rises, Consumption Need Not}
\author{LeanEconomics Org [mostly Claude]\thanks{The Lean development described here is at
\url{https://github.com/LeanEconomics/LeanEconomics}. It was written mostly by Claude (Anthropic),
an AI coding model, working under the direction of Robert Kirkby (Victoria University of
Wellington); every result reported was checked by the Lean kernel. Comments welcome.}}
\date{September 2026\\[0.5em]\small Draft}

\begin{document}
\maketitle

\begin{abstract}
\noindent
A pension that is withdrawn as the recipient's assets rise is widely said to make the household's
policy function non-monotone, and the reason usually given is that the means test makes the value
function non-concave. We show that the conclusion is right, the reason is wrong, and that which
policy misbehaves depends on the state space. With cash on hand as the state, optimal saving rises
with the state for \emph{every} continuation value, concave or not: only strict concavity of
period utility is used, and the continuation cancels. What a non-concave continuation produces is
a jump, and the jump falls on consumption, which drops discretely at the wealth where the
household gives the pension up. With assets and the earnings shock as the state, next period's
assets rise with this period's at each shock, again for every continuation; the repository's
earlier version of that statement assumed concavity, which is exactly what a means test destroys.
Monotone saving fails only when the test is assessed on assets currently held and withdraws the
pension faster than the gross return, because then cash on hand itself falls as assets rise. A
calibrated retiree confirms each case, and Australia's asset test sits on the benign side of the
threshold with a non-monotone consumption function. Every result is checked by the Lean 4 proof
assistant.
\end{abstract}

\noindent\emph{JEL:} D15, H55, C63. \quad \emph{Keywords:} means testing, monotone comparative
statics, consumption, borrowing constraints, formal verification, Lean.

\section{Introduction}

A pension that is withdrawn as the recipient's assets rise gives the household a reason to hold
fewer assets, and it is often said that this makes the household's policy function non-monotone.
The claim is usually supported by an appeal to the value function: the means test puts a kink, or
a cliff, into the reward from saving, the value function is no longer concave, and monotone
comparative statics are said to fail with it.

Three quotations fix what is actually claimed. \citet{KudrnaWoodland2011}, simulating the
Australian age pension in a general equilibrium life-cycle model, treat the geometry as settled:
``the means test for the age pension causes the budget set to be non-convex, as is well known.''
They pay for it in their solution method, whose solver ``is not guaranteed to yield a globally
optimal solution for the inter-temporal expected utility maximization problem when the budget
constraint is non-convex''; they check the solution of every household over the age of sixty by
hand to confirm that the reported optimum is global. \citet{GarstenauerSiassi2024}, modelling the
United States tax and transfer system for married couples with children, reach the same conclusion
about method: ``the presence of kinks and non-differentiabilities in the budget constraint
generated by the tax-transfer system renders Euler equation methods inapplicable.'' They solve by
discrete-state value function iteration, and because ``of the presence of these kinks and
non-convexities, it is crucial for us to capture as accurately as possible how labor supply
choices respond to taxes and transfers'' they put mothers' annual hours on a grid whose adjacent
nodes are twenty-five hours apart.

The third is closest to our question, and it cuts the other way. \citet{Wellschmied2021} studies
an asset means test in a life-cycle model and characterises the policy analytically. His Theorem 1
states that the saving policy ``is increasing'', and his proof that the value function stays
sub-differentiable uses that the policy ``is weakly increasing''. What he shows fails is
everything else. ``Despite the policy $k_h(\cdot,X)$ being monotone, the marginal propensity to
consume out of additional assets may be one for households with low income and assets''; ``the
policy function makes a jump''; and ``choices are distorted, despite satisfying first order
conditions because of the kinks in the expected value function''. A footnote records how fast this
compounds, since there is ``a rapid increase in the number of non-differentiabilities in the value
function''. The cost lands in the numerical appendix, where ``first order conditions are not
necessarily unique'', so off-grid choice is abandoned over the part of the state space the test
touches and almost nine thousand asset nodes are carried instead. His argument works through the
differential structure of the value function, following \citet{ClausenStrub2020} on how kinks
propagate under maximisation; the argument below uses no calculus at all.

The conclusion is right and the reason is wrong, and \citet{Wellschmied2021} is already halfway to
saying so. Monotonicity of the saving policy in cash on hand
survives \emph{any} continuation value whatever --- concave, kinked, discontinuous, or produced by
a means test of any shape --- because the only structure the comparison needs is strict concavity
of period utility, and the continuation cancels between the two optimality inequalities. This is
Topkis's argument \citep{Topkis1978,MilgromShannon1994} applied to the household's problem, and it
does not care what the continuation looks like.

What a non-concave continuation does produce is a jump. At the wealth where the household stops
protecting its pension and starts saving past the test, optimal saving jumps up, and because it
jumps by more than wealth has risen, consumption jumps \emph{down}. So the non-monotonicity is
real, it is in consumption, and it coexists with a perfectly monotone saving policy. We give an
exact two-period example, with the optimum at each of two wealth levels computed and verified.

Which policy misbehaves also depends on which state space one works in, and the two are not
equivalent under a means test. In the household's own state space --- assets $a$ and the earnings
shock $z$ --- next period's assets rise with this period's at each $z$, again for every
continuation. The Lean development that this paper draws on already had that statement, but with
concavity of the continuation as a hypothesis; a means test is exactly what removes that
hypothesis's premise, so we prove it without. Monotone saving in assets then fails in one case
only: when the test is assessed on the assets currently held and withdraws the pension faster than
the gross return, cash on hand itself falls as assets rise, and the policy follows it down. The
threshold is sharp, the taper rate against the gross return, and Australia's asset test sits on
the benign side of it.

Section~\ref{sec:model} sets out the retired household and the three pension schedules.
Section~\ref{sec:cash} gives the two results in the cash-on-hand state space: saving monotone
always, consumption not. Section~\ref{sec:assets} gives the result in the $(a,z)$ state space and
the exception. Section~\ref{sec:num} calibrates a retiree and checks each case. Every result
stated with a formal name has been checked by the Lean 4 proof assistant \citep{Lean4} against
Mathlib \citep{Mathlib}; the development is at
\url{https://github.com/LeanEconomics/LeanEconomics}, files
\texttt{LeanEconomics/OLG/MeansTest.lean} and \texttt{Retirement.lean}.

\section{The retired household}\label{sec:model}

A household lives $J$ periods, works until age $J_r$ and is retired after it. Period utility is
CRRA with coefficient $\mu$ and discount factor $\beta$, the gross return is $R=1+r$, and assets
may not fall below a borrowing limit, taken to be zero. In retirement labour earnings are replaced
by a pension.

\paragraph{A lump-sum pension.} If the pension is a constant $p$, the retired household faces the
same problem as the working one with earnings replaced by $p$, so the retired phase is an
income-fluctuation problem in its own right (\lean{withPension}) and the whole of the standard
theory applies to it. Two facts are worth recording. A retiree's human wealth does not depend on
the earnings shock, because the pension does not, and it never exceeds the perpetuity value of the
pension, $H_k\le p/r$ (\lean{stageHumanWealth\_withPension\_const},
\lean{stageHumanWealth\_withPension\_le}). And cash on hand is strictly increasing in assets
(\lean{withPension\_resources\_strictMono}). Both fail once the pension is means-tested.

\paragraph{Three schedules.} Write the pension as a function $p(\cdot)$ of assets and cash on hand
as
\[
  m(a)=Ra+p(a)
\]
(\lean{cash}). We consider three schedules (\lean{lumpSum}, \lean{tapered}, \lean{cliff}):
\begin{itemize}[leftmargin=1.5em]
\item a \emph{lump sum}, $p(a)=p$;
\item a \emph{taper}, $p(a)=\max\{0,\;p-\tau\max\{0,a-\bar a\}\}$, which withdraws the pension at
  rate $\tau$ per unit of assets above a free area $\bar a$ and never pays a negative amount;
\item a \emph{cliff}, $p(a)=p$ for $a\le\bar a$ and $0$ above, the limit of a taper as
  $\tau\to\infty$.
\end{itemize}

\paragraph{Where the test is assessed.} It matters whether the pension paid this period depends on
the assets held this period or on the assets carried into the next. A test on assets \emph{carried
forward} enters only the continuation value, leaving today's budget alone. A test on assets
\emph{currently held} enters today's cash on hand directly. Real asset tests are of the second
kind, assessed each period on current holdings, and the two have different consequences, as
Sections~\ref{sec:cash} and~\ref{sec:assets} show.

\paragraph{One period of the problem.} Throughout we write the household's choice as
\begin{equation}\label{eq:step}
  \max_{x\in[0,m)}\ u(m-x)+W(x),
\end{equation}
with $m$ cash on hand, $x$ next period's assets and $W$ the continuation value, and call $x$
\emph{optimal} at $m$ when it attains the maximum (\lean{IsOptimalSaving}). The feasible set
excludes $x=m$, so consumption is positive. Nothing below assumes anything about $W$ beyond what
is stated: in particular it is never assumed concave, continuous, or generated by any particular
pension schedule.

\section{Cash on hand as the state}\label{sec:cash}

\subsection{Saving is monotone, whatever the continuation}

\begin{theorem}[Monotone saving in cash on hand]\label{thm:topkis}
Let period utility be strictly concave on the positive half-line and let $W$ be \emph{any}
continuation value. If $m_1<m_2$ and $x_i$ is optimal at $m_i$ in \eqref{eq:step}, then
$x_1\le x_2$.
\formal{OLG.isOptimalSaving\_mono}
\end{theorem}

\begin{proof}[Proof sketch]
Suppose $x_2<x_1$. Both choices are feasible at both states, so adding the two optimality
inequalities,
\[
  u(m_1-x_1)+u(m_2-x_2)\ \ge\ u(m_1-x_2)+u(m_2-x_1),
\]
the continuation terms cancelling. Write $p=m_1-x_1$ and $q=m_2-x_2$. Since $x_2<x_1$ and
$m_1<m_2$ we have $p<m_1-x_2<q$, and $m_2-x_1=p+q-(m_1-x_2)$ also lies strictly between $p$ and
$q$. The two right-hand arguments are therefore strict convex combinations of $p$ and $q$ with
complementary weights, so strict concavity gives $u(m_1-x_2)+u(m_2-x_1)>u(p)+u(q)$, contradicting
the display.
\end{proof}

The proof uses nothing about $W$: it enters both optimality inequalities once and cancels. So a
value function made non-concave by a means test --- or by anything else --- cannot by itself make
saving fall with cash on hand. Note also what the hypotheses do \emph{not} include: no
differentiability, no interiority, no uniqueness of the optimum. The conclusion holds for arbitrary
optimal selections, which matters here, because at the switch point the optimum is genuinely
multi-valued.

\subsection{Consumption is not monotone}

What a non-concave continuation does produce is a jump in saving, and the jump is larger than the
increase in wealth that caused it.

\begin{theorem}[Consumption falls when wealth rises]\label{thm:cons}
There are a continuation value arising from a means-tested pension and two levels of cash on hand,
the second larger, at which optimal consumption is strictly smaller.
\formal{OLG.exists\_consumption\_decrease}
\end{theorem}

The witness is exact. Take logarithmic utility, a unit gross return, and a pension of one
withdrawn entirely above assets of one, so that next period's cash on hand is
\[
  N(x)=x+\mathbf{1}\{x\le1\},
\]
which jumps down by one as $x$ crosses the threshold, and let $W=\log N$. Then:

\begingroup\small
\begin{center}
\begin{tabular}{@{}lccc@{}}
\toprule
Cash on hand $m$ & optimal saving $x$ & consumption $m-x$ & next period's cash $N(x)$\\
\midrule
$6$ & $1$ & $5$ & $2$\\
$8$ & $4$ & $4$ & $4$\\
\bottomrule
\end{tabular}
\end{center}
\endgroup

At $m=6$ the household keeps the pension, saving the most it can without losing it; at $m=8$ the
pension is no longer worth protecting, so the household gives it up and saves the proceeds. Saving
rises from $1$ to $4$, as Theorem~\ref{thm:topkis} says it must, and since it rises by more than
wealth does, consumption falls from $5$ to $4$.

Both optima are verified rather than asserted. Because utility is logarithmic, comparing
$\log(m-x)+\log N(x)$ across $x$ is comparing the product $(m-x)N(x)$, which is a quadratic on each
side of the threshold. At $m=6$ the product is $-x^2+5x+6$ below the threshold, increasing there
and equal to $10$ at $x=1$, and $(6-x)x\le9$ above it. At $m=8$ it is $-x^2+7x+8\le14$ below and
$(8-x)x\le16$ above, attained at $x=4$. The comparisons are checked by the kernel
(\lean{isOptimalSaving\_six}, \lean{isOptimalSaving\_eight}).

\subsection{The two policies come apart}

The contrast is not an artefact of proof technique. Consumption is $m-x$ with both terms rising in
$m$, and when the continuation is non-concave the second can rise faster. Monotone consumption
therefore genuinely requires concavity of the continuation, while monotone saving does not; the
Lean development records the identity $c=m-x$ for exactly this purpose
(\lean{consumptionFnOf\_eq\_sub}) and derives consumption monotonicity only under concave slices.
So in the cash-on-hand state space the means test leaves the saving policy alone and breaks the
consumption policy, and the informal argument quoted in the introduction gets the object wrong as
well as the reason.

\section{Assets and the earnings shock as the state}\label{sec:assets}

The household's own state space is the pair $(a,z)$, assets and the earnings shock, with $z$
following a finite Markov chain and cash on hand $m(a,z)=y(z)+Ra$ in the absence of a means test
on current holdings. The question is then about $g_v(a,z)$, next period's assets.

\subsection{Next period's assets rise with this period's}

\begin{theorem}[Monotone saving in assets]\label{thm:azmono}
Let $a<a'$ both lie in the asset region. Then for every continuation $v$ and every $z$,
$g_v(a,z)\le g_v(a',z)$.
\formal{IncomeFluctuation.policyOf\_mono\_of\_lt}
\end{theorem}

This is Theorem~\ref{thm:topkis} read in the household's state space: cash on hand rises strictly
with assets, so optimal saving does. The development already contained this statement
(\lean{policyOf\_mono}) but with concave slices of the continuation as a hypothesis, and a means
test is precisely what destroys concavity, so the hypothesis had to go. The Topkis argument removes
it, at the cost of requiring the two asset levels to be distinct, which is harmless. As in
Section~\ref{sec:cash}, the continuation cancels and only strict concavity of period utility
survives.

Consumption does not inherit it, for the same reason as before: $c=m(a,z)-g_v(a,z)$ with both
terms rising in $a$, and the witness of Theorem~\ref{thm:cons} transfers directly, since at fixed
$z$ cash on hand is an increasing affine function of assets.

\subsection{When saving in assets does fail}

Theorem~\ref{thm:azmono} has one hypothesis that a means test can remove: that cash on hand rises
with assets. Where the test is assessed on assets carried forward, today's budget is untouched and
the theorem applies verbatim. Where it is assessed on assets currently held, the schedule enters
cash on hand directly, and whether it still rises is a question about the taper rate.

\begin{theorem}[A gentle taper is harmless; a cliff is not]\label{thm:taper}
If $0\le\tau<R$ then $a\mapsto Ra+p(a)$ is strictly increasing under the taper
(\lean{cash\_strictMono\_tapered}), and hence, by Theorem~\ref{thm:topkis}, saving is monotone in
assets (\lean{isOptimalSaving\_mono\_assets}). Under a cliff it is not: if $a\le\bar a<b$ and
$R(b-a)<p$ then cash on hand is strictly smaller at $b$ (\lean{cash\_lt\_cash\_of\_cliff}), and
there are two asset levels, the second larger, at which the household saves strictly less.
\formal{OLG.exists\_saving\_decrease\_in\_assets}
\end{theorem}

The threshold is sharp and it is the taper rate against the gross return. A pension withdrawn more
slowly than assets earn leaves cash on hand rising, and then no amount of damage to the value
function can make saving fall; a pension withdrawn faster makes a household with more assets poorer
in the only sense that matters for today's decision, and saving falls with it. The proof of the
first half is that the positive part is one-Lipschitz, so the pension falls by at most $\tau$ per
unit of assets while gross assets rise by $R$ (\lean{max\_zero\_sub\_le}).

The witness for the second half is again exact, and deliberately built on a \emph{concave}
continuation so that nothing is hidden in it. With a unit gross return, a pension of one withdrawn
above assets of one, and continuation $\log(x+1)$:

\begingroup\small
\begin{center}
\begin{tabular}{@{}lccc@{}}
\toprule
Assets $a$ & cash on hand $Ra+p(a)$ & optimal saving & consumption\\
\midrule
$1$ & $2$ & $1/2$ & $3/2$\\
$3/2$ & $3/2$ & $1/4$ & $5/4$\\
\bottomrule
\end{tabular}
\end{center}
\endgroup

The household with more assets has less to spend, saves less and consumes less. The
non-monotonicity here lives in the budget constraint, not in the shape of the value function, and
that is the whole difference between this case and Section~\ref{sec:cash}.

\section{A calibrated retiree}\label{sec:num}

The theorems predict exactly which of the two channels operates for which schedule, and a
calibration confirms it case by case. We solve a deterministic retiree of $25$ years by backward
induction with a global search over a grid of $6001$ asset levels --- the value function is not
concave, so a first-order condition cannot be used --- at $r=4\%$, $\beta=0.96$, relative risk
aversion $3$, a full pension of $0.40$ of mean earnings and a free area of $3$ mean-earnings units
(\texttt{numerics/meanstest.m}). The test is assessed on assets currently held, so all three
channels are live. The taper of $0.078$ per unit of assets is Australia's asset test, three dollars
a fortnight per thousand dollars.

\begingroup\small
\begin{center}
\begin{tabular}{@{}lccc@{}}
\toprule
Pension schedule & Cash on hand & Saving & Consumption \\
\midrule
Lump sum & increasing & increasing & increasing \\
Taper $0.078$ (Australia) & increasing & increasing & falls by $0.025$--$0.046$ \\
Taper $0.50$ & increasing & increasing & falls by $0.13$--$0.24$ \\
Taper $1.20$ (above $R$) & falls & falls by $0.004$ & falls by $0.55$--$0.70$ \\
Cliff & falls by $0.40$ & falls by $0.32$--$0.37$ & falls by $0.55$--$0.70$ \\
\bottomrule
\end{tabular}
\end{center}
\endgroup

Every entry is what the theorems require. Saving falls with assets in exactly the two rows where
cash on hand does, which are exactly the two rows whose taper exceeds the gross return, as
Theorem~\ref{thm:taper} says; in the other three it rises, as Theorem~\ref{thm:azmono} says it
must. Consumption falls in every row with a test of any kind, as Theorem~\ref{thm:cons} says it
can, and only in the lump-sum row is every policy monotone.

The consumption drop sits exactly at the free area and grows with age, from $0.025$ of mean
earnings at the start of retirement to $0.046$ twenty years in under the Australian taper. The
mechanism is the one the witness of Section~\ref{sec:cash} isolates: just above the free area the
effective return to saving is only $R-\tau$ until the pension is exhausted, so the household jumps
its saving towards the far side of the taper range and cuts consumption to pay for it. Under the
Australian schedule the borrowing constraint binds at about eight per cent of the state space at
every age.

The effect is not confined to retirees. A worker in the last year before retirement, facing the
Australian taper for the whole of retirement, has a monotone saving policy and a consumption policy
that falls by $0.025$ of mean earnings at the same threshold: the non-concavity is inherited
through the continuation value, and it is again consumption that carries it. That is
Theorem~\ref{thm:topkis} at work --- the continuation is as badly behaved as one likes, and saving
is monotone anyway.

\section{Conclusion}

Three statements, in the order they should be made.

Optimal saving rises with cash on hand for every continuation value, so a value function made
non-concave by a means test cannot by itself reverse a policy. The common argument from
non-concavity to non-monotonicity is therefore not valid as stated, and it is not a matter of
sharper hypotheses: the continuation cancels.

Consumption is a different object and it is genuinely non-monotone. At the wealth where the
household stops protecting its pension, saving jumps up by more than wealth has risen and
consumption drops. We give the drop exactly in a two-period example and measure it in a calibrated
retiree, where Australia's asset test produces a fall of two and a half to five per cent of mean
earnings at the free area, at every age of retirement and in the last year of work as well.

Next period's assets rise with this period's, at each earnings shock, again for every continuation.
This fails in one case only: when the test is assessed on the assets currently held and withdraws
the pension faster than assets earn, cash on hand itself falls and the policy follows. The
threshold is the taper rate against the gross return, and it is sharp.

For applied work the practical content is that a means test of the Australian kind leaves the
saving policy monotone and puts a discontinuity in the consumption function at the free area.
Solution methods that rely on monotonicity of the policy in assets are safe; methods that rely on a
continuous or monotone consumption function, or on concavity of the value function, are not. That
is close to the practice the applied literature has already settled on, which abandons Euler
equation methods and searches a fine grid \citep{Wellschmied2021,GarstenauerSiassi2024}, and the
results here say which half of that caution is needed. For
theory the content is that the two state spaces are not interchangeable once the pension depends on
assets, and that monotone comparative statics survives non-concavity far better than it is usually
given credit for.


\bibliographystyle{ecta}
\bibliography{refs}

\appendix
\section{Index of formal results}\label{app:index}

All in the namespaces \lean{LeanEconomics.OLG} and \lean{LeanEconomics.IncomeFluctuation}, files
\lean{LeanEconomics/OLG/MeansTest.lean} and \lean{Retirement.lean}. Every result depends only on
the three standard axioms (\lean{propext}, \lean{Classical.choice}, \lean{Quot.sound}) and none
uses \lean{sorry}.

\begingroup\footnotesize
\begin{longtable}{@{}p{0.38\textwidth}p{0.40\textwidth}p{0.16\textwidth}@{}}
\toprule
Result & Lean name & File \\
\midrule
\endhead
Retired household, lump-sum pension & \lean{withPension} & Retirement\\
Human wealth is shock-independent & \lean{stageHumanWealth\_withPension\_const} & Retirement\\
Perpetuity bound $H_k\le p/r$ & \lean{stageHumanWealth\_withPension\_le} & Retirement\\
Cash on hand rises with assets & \lean{withPension\_resources\_strictMono} & Retirement\\
\midrule
Pension schedules & \lean{lumpSum}, \lean{tapered}, \lean{cliff} & MeansTest\\
Cash on hand & \lean{cash} & MeansTest\\
One period of the problem & \lean{IsOptimalSaving} & MeansTest\\
Theorem~\ref{thm:topkis} & \lean{isOptimalSaving\_mono} & MeansTest\\
Monotone in assets when cash is & \lean{isOptimalSaving\_mono\_assets} & MeansTest\\
Theorem~\ref{thm:cons}, the two optima & \lean{isOptimalSaving\_six}, \lean{isOptimalSaving\_eight} & MeansTest\\
Theorem~\ref{thm:cons} & \lean{exists\_consumption\_decrease} & MeansTest\\
Consumption is cash less saving & \lean{consumptionFnOf\_eq\_sub} & MeansTest\\
Lump sum keeps cash rising & \lean{cash\_strictMono\_lumpSum} & MeansTest\\
Positive part is one-Lipschitz & \lean{max\_zero\_sub\_le} & MeansTest\\
Theorem~\ref{thm:taper}, gentle taper & \lean{cash\_strictMono\_tapered} & MeansTest\\
Theorem~\ref{thm:taper}, cliff reverses cash & \lean{cash\_lt\_cash\_of\_cliff} & MeansTest\\
Theorem~\ref{thm:taper}, the two optima & \lean{isOptimalSaving\_two}, \lean{isOptimalSaving\_threeHalves} & MeansTest\\
Theorem~\ref{thm:taper} & \lean{exists\_saving\_decrease\_in\_assets} & MeansTest\\
Theorem~\ref{thm:azmono} & \lean{policyOf\_mono\_of\_lt} & MeansTest\\
\bottomrule
\end{longtable}
\endgroup


\end{document}
